\section{Por qué las gafas inteligentes podrían convertirse en un punto de entrada}

La sección anterior explicó por qué Rokid dejó las bocinas inteligentes; esta sección explica por qué apuesta por las gafas. Zhu Mingming (祝铭明) considera que la forma de usar la IA que realmente aporta valor es en cualquier momento y lugar, fragmentada y always on. El teléfono es potente, pero sigue exigiendo un recorrido de GUI: desbloquear, abrir una app y entrar en una función. Si las gafas inteligentes se llevan puestas todo el tiempo, muchas tareas breves pueden volverse «lo que pides es lo que obtienes»: consultar la hora, recibir mensajes, escanear códigos, traducir, reconocer objetos o consultar información de contexto, todo sin sacar el teléfono.

Esto no significa que el teléfono vaya a desaparecer de inmediato, sino que la frecuencia de interacción se redistribuirá. El teléfono podría conservar la comunicación, las tareas pesadas, la entrada compleja y las operaciones en pantalla grande; las gafas absorberían primero gran parte de las interacciones fragmentadas de menos de cinco minutos. El juicio de Zhu Mingming es este: para quienes ya usan lentes, las gafas con IA podrían volverse algo de uso generalizado en un plazo de tres a cinco años; a quienes no quieren usar lentes no hay que intentar convencerlos a la fuerza, sino esperar a que la diferencia de capacidades sea lo bastante evidente y dejar que ellos decidan.

\begin{figure}[H]
\centering
\includegraphics[width=0.96\textwidth]{smart-glasses-stage.png}
\caption{Etapa de las gafas inteligentes: entre BlackBerry y el iPhone 1, el punto de entrada está tomando forma. Diagrama conceptual propio, elaborado a partir de la conversación en 01:05:45--01:13:05.}
\end{figure}

\begin{knowledgebox}{Lectura de la figura: que la etapa sea inmadura no significa que la dirección sea errónea}
La figura va de Feature a BlackBerry, iPhone 1, Platform y Mass Market, y expresa que las gafas inteligentes todavía están en la etapa temprana en la que se forma un punto de entrada. Los productos actuales no son necesariamente maduros, pero si se confirman el tiempo de uso, los desarrolladores y la migración de tareas fragmentadas, tienen la oportunidad de pasar de hardware de función a plataforma.
\end{knowledgebox}

\subsection{Acceso directo a la información: reducir al mínimo el costo del sistema}

Esta sección descompone la «comodidad» en costos de interacción analizables. Zhu Mingming pone como ejemplos escanear un código, consultar el clima, ver la hora y leer mensajes: en el teléfono, el recorrido suele incluir sacar el dispositivo, desbloquearlo, encontrar la app, entrar en la función y confirmar la operación; con unas gafas que se llevan puestas todo el tiempo, basta con invocarlas por voz, confirmar visualmente o tocarlas ligeramente para ver el resultado. Unos segundos de diferencia importan poco en una sola tarea, pero cuando todas las tareas son acciones fragmentadas de menos de cinco minutos, esos segundos se acumulan hasta provocar la migración del punto de entrada.

\[
T_{\text{task}} = T_{\text{wake}} + T_{\text{path}} + T_{\text{input}} + T_{\text{output}} + T_{\text{confirm}}
\]

Aquí, $T_{\text{task}}$ es el tiempo total para completar una tarea breve; $T_{\text{wake}}$ es el tiempo para sacar o activar el dispositivo; $T_{\text{path}}$ es el costo del recorrido desde el punto de entrada del sistema hasta la función concreta; $T_{\text{input}}$ es el tiempo para expresar la intención; $T_{\text{output}}$ es el tiempo para recibir el resultado; $T_{\text{confirm}}$ es el tiempo para confirmar y corregir errores. Lo esencial de la vía de las gafas no es llevar cada término a cero, sino reducir mucho $T_{\text{wake}}$ y $T_{\text{path}}$, y usar la pantalla para disminuir la incertidumbre de $T_{\text{output}}$ y $T_{\text{confirm}}$.

\begin{importantbox}{El punto de entrada no depende de la eficiencia de una sola vez, sino de la frecuencia multiplicada por la eficiencia}
Si una acción ocurre una sola vez al día, ahorrar cinco segundos no tiene importancia; si las notificaciones, la hora, el escaneo de códigos, la traducción, el reconocimiento y la búsqueda ocurren decenas de veces al día, el costo del recorrido se convierte en un criterio para elegir el punto de entrada. El juicio de Rokid es que las gafas inteligentes primero sustituirán una gran cantidad de tareas ligeras y fragmentadas, no las tareas pesadas del teléfono de inmediato.
\end{importantbox}

\subsection{Por qué importa la pantalla}

Una de las diferencias centrales en la definición de producto de las gafas inteligentes entre China y Estados Unidos es si necesitan pantalla. Rokid considera que, si las gafas son un punto de entrada a la IA, la pantalla es indispensable. Tomemos la traducción: sin pantalla, el sistema debe esperar a tener suficiente confianza antes de leer la traducción en voz alta, y en la conversación aparecen silencios de varios segundos hasta diez segundos; con pantalla, puede mostrar mientras escucha y corregir sobre la marcha, y el usuario tolera resultados intermedios que al principio no son del todo precisos. La pantalla no está para lucirse, sino para que la salida de la IA pase de ser solo voz a una interfaz que se puede leer de un vistazo, corregir y usar en paralelo.

\begin{knowledgebox}{Términos clave: AUI y GUI}
\noindent\begin{tabularx}{\textwidth}{p{0.18\textwidth}p{0.34\textwidth}X}
\toprule
Modo de interacción & Ventajas & Limitaciones \\
\midrule
GUI & Visible, editable y fácil de leer de un vistazo; adecuada para tareas complejas & Requiere recorridos y operación de la interfaz; el costo del sistema es alto. \\
AUI & Acceso directo por voz; adecuada para tareas breves y uso sin manos & La salida es secuencial; no es adecuada para contenido largo ni para resultados intermedios de baja confianza. \\
Pantalla en gafas + IA & Combina entrada por voz y salida visual & Exige que la comodidad al llevarlas puestas, la autonomía de la batería, la pantalla y la privacidad superen juntas el umbral. \\
\bottomrule
\end{tabularx}
\end{knowledgebox}

\begin{warningbox}{No entender la pantalla como «una pantalla más»}
El verdadero valor de la pantalla no es trasladar la del teléfono frente a los ojos, sino cambiar la estructura temporal de la salida de la IA. La salida por voz tiene que ser secuencial, completa y de confianza relativamente alta; la salida visual puede mostrarse de forma gradual, admitir correcciones intermedias y permitir que el usuario la lea de un vistazo. La traducción es un caso representativo porque expone a la vez la latencia, la confianza y la posibilidad de corregir errores.
\end{warningbox}

\noindent\begin{tabularx}{\textwidth}{p{0.18\textwidth}p{0.34\textwidth}X}
\toprule
Escenario & Recorrido sin pantalla & Recorrido con pantalla \\
\midrule
Traducción & Se lee en voz alta solo cuando el modelo tiene suficiente confianza; la conversación entre ambas partes se interrumpe & Puede mostrar primero una traducción aproximada y corregirla según el contexto; el usuario la lee de un vistazo. \\
Notificaciones & El aviso por voz interrumpe la situación, y revisar el teléfono alarga el recorrido & Aviso breve frente a los ojos; el usuario puede ignorarlo o atenderlo. \\
Escaneo de códigos/pago & Hay que abrir el teléfono, buscar el punto de entrada y confirmar & Se dirige la mirada al código QR y se confirma por voz o con un toque ligero. \\
Reconocimiento de objetos & Tomar una foto con el teléfono o abrir la cámara & Las gafas forman la consulta directamente a partir de la dirección de la mirada. \\
\bottomrule
\end{tabularx}

\subsection{Límites de uso: no intentar convencer a quien no quiere usar lentes}

Esta sección aborda una objeción realista: mucha gente no quiere usar lentes, e incluso quienes tienen miopía optan por lentes de contacto o cirugía. La respuesta de Zhu Mingming es mesurada: en la etapa actual no hay que intentar convencer a la fuerza a ese grupo. Primero se atiende a quienes ya aceptan lentes con armazón, sumando capacidades de IA a un hábito que ya tienen; cuando las personas de su entorno que usan lentes obtengan una ventaja informativa evidente, quienes no quieren usarlos podrán volver a decidir. La estrategia de producto no es «que todos los usen de inmediato», sino empezar por el grupo para el que resulta más natural, los usuarios nativos.

\begin{knowledgebox}{Experiencia práctica: la curva de adopción empieza por el grupo con menos resistencia}
Los primeros usuarios de un producto de hardware no son necesariamente el grupo más grande, sino el que opone menos resistencia, da la retroalimentación más intensa y está dispuesto a tolerar las imperfecciones iniciales. Para las gafas inteligentes, quienes están dispuestos a usar lentes con armazón, los entusiastas de la tecnología y los usuarios frecuentes de viajes de trabajo, traducción o reuniones son mejores usuarios principales tempranos que quienes rechazan por completo los lentes.
\end{knowledgebox}

\noindent\begin{tabularx}{\textwidth}{p{0.20\textwidth}p{0.34\textwidth}X}
\toprule
Tipo de usuario & Resistencia inicial & Estrategia de producto \\
\midrule
Quienes ya usan lentes & La más baja; solo hay que convertir el uso existente en capacidad de IA & Priorizar comodidad, graduación, peso y uso diario. \\
Quienes están dispuestos a usarlos en algunos escenarios & Media; los escenarios clave deben ser lo bastante sólidos & Escenarios de alto valor como viajes, traducción, reuniones, navegación y escaneo de códigos. \\
Quienes no quieren usar lentes & La más alta; hay resistencia tanto fisiológica como psicológica & No forzar su convencimiento; ampliar el alcance cuando la diferencia de capacidades y el formato hayan madurado. \\
\bottomrule
\end{tabularx}

\subsection{Diferencias de definición de producto entre China y Estados Unidos}

Muchas gafas inteligentes estadounidenses siguen más bien la lógica de las sunglasses: exteriores, redes sociales, fotos y videos para compartir; no necesariamente son always on ni tratan la pantalla como indispensable. Rokid se inclina más por la lógica china de los lentes: muchos usuarios ya usan lentes, y llevarlos puestos durante largos periodos, en interiores y exteriores, resulta más natural; por eso, desde el primer día, la definición de producto se orienta a IA más AR, con énfasis en la pantalla, el asistente y un punto de entrada a la información que se lleva encima. La vía de Meta/Ray-Ban es una herramienta de producción de contenido social; la de Rokid es el consumo personal de información y el asistente de IA.

\begin{figure}[H]
\centering
\includegraphics[width=0.94\textwidth]{china-us-glasses-definition.png}
\caption{Diferencias en la definición de las gafas inteligentes entre China y Estados Unidos: difieren la definición de producto, los canales y los escenarios de uso. Diagrama conceptual propio, elaborado a partir de la conversación en 01:23:38--01:41:35.}
\end{figure}

\begin{knowledgebox}{Lectura de la figura: que ambas parezcan lentes no significa que sean la misma especie}
A la izquierda, la vía china enfatiza a la población con miopía, el uso durante todo el día, la pantalla y el asistente de IA; a la derecha, la vía estadounidense entra con más facilidad desde las sunglasses, la grabación en redes sociales y la moda de marca. Ambas podrían converger con el tiempo, pero en la etapa temprana difieren en la definición de producto, los escenarios de uso y la prioridad de las funciones.
\end{knowledgebox}

\begin{importantbox}{La definición de producto determina la pila tecnológica}
Si el producto es una herramienta de grabación para redes sociales, lo central son la cámara, la marca, la cadena para compartir y los avisos de privacidad; si el producto es un punto de entrada a la IA, lo central son la pantalla, el sistema operativo, la retroalimentación de baja latencia, el consumo de energía, la comodidad al llevarlas puestas y la integración de modelos. Aunque ambas se llamen gafas inteligentes, sus prioridades técnicas pueden ser completamente distintas.
\end{importantbox}

\subsection{Resumen de la sección}

La lógica de punto de entrada de las gafas inteligentes se basa en estar always-on, en el acceso directo a la información y en la migración de tareas fragmentadas. La capacidad de pantalla determina si la salida de la IA puede pasar de voz secuencial a una interacción visible, corregible y capaz de convertirse en plataforma; los límites de uso, a su vez, determinan que el grupo inicial no puede ampliarse sin límite y que hay que empezar por los usuarios con menos resistencia.
